\documentclass{article}
\usepackage{amsmath}
\usepackage{xcolor}
\usepackage{hyperref}

% Switch to XeLaTeX to print Unicode
\usepackage{fontspec}
\setmainfont{Latin Modern Roman}
\newfontfamily\cjkfont{Noto Sans CJK SC}
\newcommand{\cjk}[1]{{\cjkfont #1}}

% Size and space
\usepackage{titlesec}
\titleformat{\subsection}{\small\bfseries}{\thesubsection}{1em}{}
\titlespacing*{\subsection}{0pt}{24pt}{12pt}

\begin{document}

\subsection*{1. Wie viele Zeichen sind in UTF-8 darstellbar? Wieso ist das in der Praxis weniger als theoretisch möglich?}

UTF-8 erlaubt Codepoint-Sequenzen von 1 bis 4 Bytes.
Der größte theoretisch darstellbare Wert wäre alle 4 Bytes mit Einsen aufzufüllen:

\begin{figure}[h]
\centering
\texttt{%
\textcolor{red}{11110}\textcolor{blue}{xxx} 
\textcolor{red}{10}\textcolor{blue}{xxxxxx} 
\textcolor{red}{10}\textcolor{blue}{xxxxxx} 
\textcolor{red}{10}\textcolor{blue}{xxxxxx}
}
\caption{\textcolor{red}{Kontrollbits}, \textcolor{blue}{Wertbits}}
\end{figure}

21 Wertbits ergeben $2^{21} - 1$, also ca. 2,1 Millionen.

In der Tat sind laut Unicode-Standard etwa 1,1 Millionen Zeichen darstellbar.

Wieso ist das so? 
 
\textbf{Gründe für die Reduzierung:}
Die kurze Antwort lautet: weil es die Unicode Spezifikation so definiert\footnote{\url{https://www.unicode.org/versions/Unicode16.0.0/core-spec/chapter-3/#G740}}: 

\begin{quote}
Each encoding form maps the Unicode code points U+0000..U+D7FF and U+E000..U+10FFFF to unique code unit sequences. 
\end{quote}
\textit{(The Unicode Standard, Version 16.0, Chapter 3)}

Primärquellen für diese Entscheidung sind nur schwer aufzufinden. Sekundärquellen scheinen sich einig darüber, dass die Anzahl an darstellbaren Zeichen in UTF-8 reduziert wurde, um das Encoding mit UTF-16 kompatibel zu machen. UTF-16 hat einen beträchtlich kleineren Spielraum. Wären gewisse Zeichen nur durch UTF-8 aber nicht durch UTF-16 darstellbar, würde die Komplexität des Systems in die Höhe schießen und das Ziel einer einheitlichen Darstellung von Zeichen wäre misslungen.

Außerdem repräsentiert die Obergrenze von 10FFFF eine Sicherheitsmaßnahme. Die UTF-8 Spezifikation (RFC 3629) begründet\footnote{\url{https://datatracker.ietf.org/doc/html/rfc3629#section-1}}:

\begin{quote}
 Another security issue occurs when encoding to UTF-8: the ISO/IEC
   10646 description of UTF-8 allows encoding character numbers up to
   U+7FFFFFFF, yielding sequences of up to 6 bytes.  There is therefore
   a risk of buffer overflow if the range of character numbers is not
   explicitly limited to U+10FFFF or if buffer sizing doesn't take into
   account the possibility of 5- and 6-byte sequences
\end{quote}

Weitere Reduktion ergibt sich aus der Auslassung von Codepoint-Sequenzen U+D800..U+DFFF. Dies hat wiederum damit zu tun, dass dieses Intervall für sog. Surrogate in UTF-16 reserviert sind. (Surrogate sind in UTF-16 zwei 16-Bit-Einheiten, die es UTF-16 ermöglichen, Codepoints jenseits von FFFF zu kodieren.) Insgesamt werden 2048 gesperrt. Damit dominiert die Obergrenze bei Weitem.

Weitere Quellen:

- Ein Brief der Unicode Gruppe an das ISO Konsortium, in dem sie bitten, die Range der Code Positions auf 10FFFF zu reduzieren\footnote{\url{https://www.unicode.org/L1/L2000/00079-n2175.html}}  

\subsection*{2. Bytestream}

Aufgabe:
\begin{quote}
Gegeben sei folgender Ausschnitt eines Bytestreams: 41 C3 84 E4 B8 AD
C3 A9 CE A9. Bestimmen Sie für jedes Byte, ob es ein Startbyte oder ein
Folgebyte ist. Folgern Sie, wie viele Zeichen die Nachricht enthält
\end{quote}

\noindent Lösung:

Wir stellen die Bytes in Binärform dar.

Startbytes beginnen mit \texttt{0, 10, 110, 1110}.

\begin{align*}
\texttt{41} &= 4 * 16 + 1 &= \texttt{0100 0001} & \quad \text{Startbyte} \\
\texttt{C3} &= 12 * 16 + 3 &= \texttt{1100 0011} & \quad \text{Startbyte} \\
\texttt{84} &= 8 * 16 + 4 &= \texttt{1000 0100} & \quad \textcolor{gray}{\mathrm{Folgebyte}} \\
\texttt{E4} &= 14 * 16 + 4 &= \texttt{1110 0100} & \quad \text{Startbyte} \\
\texttt{B8} &= 11 * 16 + 8 &= \texttt{1011 1000} & \quad \textcolor{gray}{\mathrm{Folgebyte}} \\
\texttt{AD} &= 10 * 16 + 13 &= \texttt{1010 1101} & \quad \textcolor{gray}{\mathrm{Folgebyte}} \\
\texttt{C3} &= 12 * 16 + 3 &= \texttt{1100 0011} & \quad \text{Startbyte} \\
\texttt{A9} &= 10 * 16 + 9 &= \texttt{1010 1001} & \quad \textcolor{gray}{\mathrm{Folgebyte}} \\
\texttt{CE} &= 12 * 16 + 14 &= \texttt{1100 1110} & \quad \text{Startbyte} \\
\texttt{A9} &= 10 * 16 + 9 &= \texttt{1010 1001} & \quad \textcolor{gray}{\mathrm{Folgebyte}} \\
\end{align*}

5 Startbytes entspricht 5 Zeichen:

\begin{align*}
    \texttt{41} \\
    \texttt{C384} \\
    \texttt{E4B8AD} \\
    \texttt{C3A9} \\
    \texttt{CEA9} \\
\end{align*}

Das ergibt die Nachricht AÄ\cjk{中}éΩ.

\subsection*{3. Selbstsynchronisation}

Aufgabe:

\begin{quote}
    UTF-8 ist ein selbstsynchronisierendes Kodierungsformat, das bedeutet es ist
möglich, beim Lesen eines Bytestroms an einer beliebigen Stelle einzusteigen
und die Grenzen der Zeichen zu erkennen. Erläutern Sie, wie dies möglich
ist.
\end{quote}

\noindent Lösung:

Wie wir in der vorigen Aufgabe gesehen haben, kann ein Programm an jedem Byte ablesen, ob es sich um ein Startbyte oder Folgebyte handelt. Wenn das Programm in der Mitte eines Zeichens landet, kann es einfach den Stream weiterlesen, bis es ein Byte findet, das mit einer Startsequenz beginnt und von da an parsen.

\subsection*{4. Bytegrenzen}

Aufgabe:

\begin{quote}
Erklären Sie, warum dasselbe nicht immer möglich ist, wenn man ohne
Kenntnis der Bytegrenzen bei einem zufälligen Bit beginnt zu lesen.
\end{quote}

\noindent Lösung:

Das UTF-8 Schema basiert darauf, dass die Kontrollbits immer am Anfang eines Bytes stehen. 
Wenn wir nicht wissen, wo ein Byte anfängt, bricht das System zusammen: die Applikation
kann ein Bit nicht mehr eindeutig als Kontrollbit identifizieren. Beispiel: Die Byte-Sequenz \texttt{11000011 10000100} (Ä) wird zu \texttt{10000111 0000100x}... wenn man bei Bit 2 startet - die Kontrollmuster des Originals sind nicht mehr erkennbar.

\end{document}